\chapter{Anforderungen und Analyse}

\section{User Story}

Aus dem Dokument \textit{Aufgabenbeschreibung} und den Gesprächen mit Prof. Dr. Tholen ergab sich folgende User Story:

\begin{center}
\begin{tcolorbox}[colback=gray!10, colframe=black, width=0.9\textwidth, sharp corners=south]
\textbf{As a professor of Applied Artificial Intelligence,}\\
I want a \textbf{graphical user interface} that demonstrates \textbf{explainable AI (xAI),}\\
so I can visually teach how AI models explain their decisions to my students.
\end{tcolorbox}
\end{center}

Die Hauptstakeholders dieser User Story sind somit:
\begin{enumerate}
    \item ein Professor der Angewandten Künstlichen Intelligenz, der das Tool in der Lehre einsetzen möchte
    \item Studierende, die das Tool zur Unterstützung ihres Lernprozesses verwenden
\end{enumerate}
    
Durch eine interaktive grafische Benutzeroberfläche (GUI) soll das Verständnis für Explainable AI (XAI) gefördert werden, indem die Funktionsweise und Ergebnisse von XAI-Verfahren auf Black-Box-Modellen visuell dargestellt werden.\\


\section{Anforderungsanalyse}

Die Anforderungen an das interaktive Tool wurden erhoben, mithilfe der MoSCoW–Technik priorisiert und durch den Stakeholder bestätigt.

In dieser Hinsicht wurden die Anforderungen in folgende Kategorien eingeteilt:

\subsection{Must-Have Anforderungen}
\todox[]{Anforderungen einfügen und begründen}

\subsection{Should-Have Anforderungen}
\todox[]{Anforderungen einfügen und begründen}
\subsection{Could-Have Anforderungen}
\todox[]{Anforderungen einfügen und begründen}
\subsection{Won't-Have Anforderungen}
\todox[]{Anforderungen einfügen und begründen}

